% TotalSegmentator -> start points -> in-house tracer (Kit Mills Bransby). Not the author's
% contribution, so only outputs and how it was run on CGPS; the author's check is sec:method-medis.
% Each red [Qn] blank is answered by Qn in the list below; fill the blanks, then delete the list.
% Red notes need \usepackage{xcolor}. Remove before submission.

\newpage

\providecommand{\clarify}[1]{\textcolor{red}{\textbf{Clarification needed:} #1}}
\providecommand{\ask}[1]{\textcolor{red}{[#1]}}
\section{Coronary Centerline and Lumen Extraction}
\label{sec:method-tracer}

To extract the coronary centerlines and lumen from the baseline CCTA volumes, an in-house tracer was applied to all scans of the Herlev--{\O}sterbro cohort (Section~\ref{sec:data-cohort}). The in-house tracer was developed and run by Kit Mills Bransby. The method is described here so that the full pipeline is shown. 


\clarify{Questions for Kit:
\begin{enumerate}
  \item Can you give a brief description of the tracer. Should we call it in-house tracer or do you have a better name for it?
\end{enumerate}}
